\section{Calculus as a stable closure}
\label{sec:exhibit-calculus}

This section applies the engine to the most familiar case: how differentiation and integration arise as stable closures of discrete protocols.
It has three parts: an exact ledger for products (Sections~\ref{subsec:ledger}--\ref{subsec:poly}), a numerical test of whether that ledger selects derivative-like operators from a population of candidates (Section~\ref{subsec:stencils}), and the corresponding audit for sums (Section~\ref{subsec:integration}).

\subsection{Discrete substrate}
Fix a step $h\ne 0$ and functions on a set where translation by $h$ makes sense (a grid, a line, or a field).
The substrate provides the shift $(T_hf)(x)=f(x+h)$, the difference $\Delta_h=T_h-I$, and the scaled difference
\[
  \delta_hf(x)\;:=\;\frac{f(x+h)-f(x)}{h}.
\]
The factor $h^{-1}$ is the simplest operator rewrite (\Pone): without it, the difference of a smooth function vanishes as $h\to 0$ and nothing survives refinement.
Refinement $h\to 0$ is the staging (\Pfour).

\subsection{The ledger: an exact product rule with remainder}
\label{subsec:ledger}
At finite $h$ the scaled difference is not a derivation, but its failure is exactly computable.

\begin{lemma}[Finite-difference Leibniz identity]\label{lem:finite-diff-leibniz}
Let $F$ be a field, $h\in F$ with $h\ne 0$, and $f,g:F\to F$. For all $x\in F$,
\[
  \delta_h(fg)(x)\;=\;\delta_hf(x)\,g(x)\;+\;f(x)\,\delta_hg(x)\;+\;h\,\delta_hf(x)\,\delta_hg(x).
\]
Moreover, $\delta_h$ is linear and annihilates constants.
\end{lemma}

\begin{proof}
Expand $f(x+h)g(x+h)-f(x)g(x)=\bigl(f(x+h)-f(x)\bigr)g(x)+f(x)\bigl(g(x+h)-g(x)\bigr)+\bigl(f(x+h)-f(x)\bigr)\bigl(g(x+h)-g(x)\bigr)$ and divide by $h$.
\end{proof}

\noindent The Lean statements are \leanname{Diff.delta_mul_leibniz}, \leanname{Diff.delta_add}, \leanname{Diff.delta_const_mul}, and \leanname{Diff.delta_const}, in \path{lean/Diff/FiniteDifference.lean}.
The identity is purely algebraic; it needs neither limits nor characteristic zero.

The remainder $h\,\delta_hf\,\delta_hg$ is the ledger term (\Psix) for the passage from differences to derivatives.
It carries an explicit factor of the stage parameter, so it vanishes as $h\to 0$ whenever the two factor quotients stay bounded.
That observation is the heart of the closure story, but turning it into a theorem requires care about what ``the limit'' is.

\subsection{From a vanishing remainder to derivations}
\label{subsec:bridge}
The first route identifies the packaged operator with ordinary pointwise limits.

\begin{proposition}[Limit bridge]\label{prop:limit-bridge}
Let $l$ be a filter on an index set, and let $h_i\in\mathbb R$ with $h_i\to 0$ along $l$ and $h_i\ne 0$ eventually.
Let $f,g:\mathbb R\to\mathbb R$ and $x\in\mathbb R$, and suppose $\delta_{h_i}f(x)\to a$ and $\delta_{h_i}g(x)\to b$. Then
\[
  \delta_{h_i}(fg)(x)\;\longrightarrow\;a\,g(x)+f(x)\,b .
\]
If, in addition, $l$ is nontrivial and $\delta_{h_i}(fg)(x)\to c$, then $c=a\,g(x)+f(x)\,b$.
\end{proposition}

\begin{proof}
By Lemma~\ref{lem:finite-diff-leibniz}, eventually
\[
  \delta_{h_i}(fg)(x)=\delta_{h_i}f(x)\,g(x)+f(x)\,\delta_{h_i}g(x)+h_i\,\delta_{h_i}f(x)\,\delta_{h_i}g(x).
\]
The first two terms converge to $a\,g(x)$ and $f(x)\,b$, and the third to $0\cdot a\cdot b=0$, by continuity of addition and multiplication. Uniqueness of limits in a Hausdorff space along a nontrivial filter gives the second claim.
\end{proof}

\noindent These are \leanname{Diff.tendsto_delta_mul} and \leanname{Diff.leibniz_of_tendsto_delta}.
Two features matter.
The convergence of the product quotient is \emph{derived}, not assumed.
And the nontriviality hypothesis is needed for the identification: along the trivial filter every value is a limit.

Combined with the linearity in Lemma~\ref{lem:finite-diff-leibniz}, Proposition~\ref{prop:limit-bridge} shows that on any algebra of functions where the limits $D f(x)=\lim\delta_{h_i}f(x)$ exist, all taken along the same nontrivial filter, $D$ is linear, kills constants, and satisfies $D(fg)=(Df)g+f(Dg)$.
The values of $D$ may lie in a different regularity class from its inputs; differentiation maps $C^1$ to $C^0$, while polynomials are closed under formal differentiation.
A limit along one sequence of steps is also weaker than a derivative along all steps; the filter must be stated.

Bounded factor quotients make the ledger term vanish, but they do not make the quotients converge.

\begin{example}[A vanishing ledger without a derivative]\label{ex:bounded-not-convergent}
Let $f(0)=0$ and $f(x)=x\sin(\log|x|)$ for $x\ne 0$. Then $\delta_hf(0)=\sin(\log|h|)$ is bounded by $1$ but has no limit as $h\to 0$. With $g=f$, the ledger term $h\sin^2(\log|h|)$ tends to $0$, yet $f'(0)$ does not exist.
\end{example}

So a shrinking ledger alone does not construct a derivative.
It does, however, construct a derivation if the target is chosen to be the space where bounded families live.

\begin{proposition}[A quotient-valued derivation]\label{prop:quotient-derivation}
Let $A$ be the real algebra of bounded Lipschitz functions $\mathbb R\to\mathbb R$ and $B$ the Banach algebra of bounded functions with the uniform norm. Let
\[
  Q\;=\;\ell^\infty(\mathbb N;B)\,/\,c_0(\mathbb N;B),
\]
regarded as an $A$-module through pointwise multiplication by $A$ on each term (which preserves $c_0$ because elements of $A$ are bounded).
Fix $h_n\ne 0$ with $h_n\to 0$ and set $D(f)=[(\delta_{h_n}f)_n]\in Q$. Then $D:A\to Q$ is a well-defined $\mathbb R$-linear derivation:
\[
  D(fg)\;=\;D(f)\,g+f\,D(g)\qquad\text{in }Q.
\]
\end{proposition}

\begin{proof}
For $f\in A$, $\|\delta_{h_n}f\|_\infty\le\Lip(f)$, so the sequence lies in $\ell^\infty(\mathbb N;B)$; since $A$ is an algebra, the same holds for $fg$. Linearity follows from the exact linearity of $\delta_h$. By Lemma~\ref{lem:finite-diff-leibniz}, the difference between $\delta_{h_n}(fg)$ and $(\delta_{h_n}f)g+f(\delta_{h_n}g)$ is $h_n\,\delta_{h_n}f\,\delta_{h_n}g$, whose uniform norm is at most $|h_n|\Lip(f)\Lip(g)\to 0$. So the difference lies in $c_0(\mathbb N;B)$ and vanishes in $Q$.
\end{proof}

This is the precise sense in which ``derivations are what survives when the multiplicative ledger vanishes'': the derivation identity holds exactly, but in a quotient that does not pretend to be $\mathbb R$-valued and that need not consist of classical derivatives.
The construction is on fixed functions in $A$; it makes no claim about descent for arbitrary families of inputs (Example~\ref{ex:no-descent}).
Proposition~\ref{prop:quotient-derivation} is proved here and is not mechanized; its pointwise counterpart, Proposition~\ref{prop:limit-bridge}, is.

\subsection{An algebraic anchor: derivations of polynomial rings}
\label{subsec:poly}
In a purely algebraic setting, imposing the product rule leaves remarkably little freedom.

\begin{theorem}[Derivations of $R\lbrack X\rbrack$]\label{thm:derivation-determined-by-X}
Let $R$ be a commutative semiring, and let $\mathrm{Der}_R(R[X],R[X])$ denote the $R$-derivations of $R[X]$, that is, $R$-linear maps that satisfy the product rule and annihilate the constant polynomials ($D(1)=0$, which over a general semiring does not follow from the other two conditions). Evaluation at the indeterminate,
\[
  \mathrm{Der}_R\bigl(R[X],R[X]\bigr)\;\xrightarrow{\ \sim\ }\;R[X],\qquad D\longmapsto D(X),
\]
is an $R$-linear equivalence, with inverse $p\mapsto\bigl(q\mapsto p\,q'\bigr)$, where $q'$ is the formal derivative. In particular, a derivation is determined by $D(X)$, and the derivation with $D(X)=1$ is the formal derivative.
\end{theorem}

\begin{proof}[Proof sketch]
A derivation is $R$-linear and kills scalars. Iterating the product rule gives $D(X^n)=nX^{n-1}D(X)$, so by linearity $D(q)=q'\,D(X)$ for every polynomial $q$. Conversely, $q\mapsto p\,q'$ is a derivation for every $p$. No division and no assumption on the characteristic is needed.
\end{proof}

\noindent The Lean counterparts, in \path{lean/Derivations/Polynomial.lean}, are:
\begin{itemize}
  \item uniqueness: \leanname{Derivations.derivation_ext_X};
  \item the equivalence, imported from mathlib's proved equivalence:\\ \leanname{Derivations.polynomialDerivationEquiv};
  \item its evaluation lemmas:\\ \leanname{Derivations.polynomialDerivationEquiv_apply_X},\\ \leanname{Derivations.polynomialDerivationEquiv_symm_apply_X};
  \item the inverse formula:\\ \leanname{Derivations.polynomialDerivationEquiv_symm_apply};
  \item the normalized statement:\\ \leanname{Derivations.derivation_eq_derivative_of_X_eq_one}.
\end{itemize}
Classically this is the computation $\Omega_{R[X]/R}\cong R[X]\,dX$ of K\"ahler differentials \cite{StacksProject00RM}.

The theorem expresses rigidity of an algebraic family of operators, not finite-dimensionality ($R[X]$ need not be finite-dimensional over $R$), and it does not classify operators on arbitrary real functions.
In the analytic setting the same rigidity is mediated by the ledger and the packaging contract, which is what the next experiment probes.

\subsection{Does the Leibniz ledger select derivatives? A matched stencil experiment}
\label{subsec:stencils}
We now ask a falsifiable version of the closure claim: among many local linear operators that are stable under refinement, does requiring a shrinking Leibniz defect single out the first-derivative ones?

\paragraph{Candidates.}
A candidate is a stencil of half-width $m=\StencilMatchedM$,
\[
  (L_hf)_i\;=\;h^{-k}\sum_{j=-m}^{m}c_j\,f_{i+j},
\]
applied at interior grid points only, so no boundary rule is involved.
Each candidate has a \emph{declared} scaling order $k\in\{0,1,2\}$, fixed when it is generated, and its coefficients satisfy the moment conditions $M_p:=\sum_jj^pc_j=k!\,\delta_{pk}$ for $p=0,1,2,3$.
For a $C^4$ function $f$, Taylor's theorem then gives $L_hf(x)=f^{(k)}(x)+O(h^{4-k})$, with a constant proportional to $\|f^{(4)}\|_\infty\sum_j|c_j|\,|j|^4$; so each candidate is consistent with $f$, $f'$, or $f''$ according to its order.
We draw $\StencilMatchedPerOrder$ candidates of each order by projecting Gaussian coefficient vectors onto these moment conditions (seed $0$), giving a fixed population of $\StencilMatchedTotal$ stencils.

\paragraph{Gates.}
Each candidate is applied on uniform grids $x_i=ih$ in $[0,1)$ with $N_0=\StencilMatchedNZero$, $2N_0$, and $4N_0$ points, to the test family $\{1,\,x,\,x^2,\,\sin 2\pi x,\,e^x\}$, with outputs compared at aligned interior centers.
For consecutive grids, the refinement mismatch is $\mathrm{RMS}(\text{coarse}-\text{aligned fine})/(\mathrm{RMS}(\text{coarse})+10^{-8})$, maximized over the family; write $E_{01}$ and $E_{12}$ for its values on the first and second pair of grids. The \emph{stability gate} requires $E_{12}<E_{01}$ (or $E_{12}<10^{-8}$) and $E_{12}<0.25$, and it requires a nontrivial output: the RMS of the pooled finest-grid outputs must exceed $10^{-6}$.
The \emph{Leibniz gate} adds a product-rule audit.
For the pairs $(x,x^2)$, $(\sin 2\pi x,e^x)$, $(x^2,e^x)$, and $(1,e^x)$ we compute
\[
  \Def_{\mathrm L}\;=\;\frac{\mathrm{RMS}\bigl(L(fg)-(Lf)g-f(Lg)\bigr)}{\mathrm{RMS}\bigl(L(fg)\bigr)+\mathrm{RMS}\bigl((Lf)g\bigr)+\mathrm{RMS}\bigl(f(Lg)\bigr)+10^{-12}},
\]
maximized over the pairs; normalizing by all three terms avoids dividing by values near stationary points.
Writing $\Def_0,\Def_1,\Def_2$ for its values on the three grids, the gate requires $\Def_2<\StencilMatchedDTwoMax$, $\Def_2<\StencilMatchedRatioMax\,\Def_1$, and $\Def_1<\StencilMatchedRatioMax\,\Def_0$.
After a candidate is accepted or rejected, we record which of $f$, $f'$, $f''$ its finest-grid output fits best up to a scalar.
This fit label is an output only; it never enters a gate.

\paragraph{Result.}
Table~\ref{tab:stencil-matched} and Figure~\ref{fig:stencil} show the outcome.
The stability gate alone keeps $\StencilMatchedStableTotal$ of the $\StencilMatchedTotal$ candidates, spread over all three orders ($\StencilMatchedStableZero$, $\StencilMatchedStableOne$, and $\StencilMatchedStableTwo$).
Adding the Leibniz gate keeps $\StencilMatchedLeibnizTotal$: all of the first-order stencils and none of the others.
The best-fit label agrees with the declared order for all $\StencilMatchedLabelAgreement$ candidates.

% Auto-generated by scripts/export_results_tex.py from notes/math_review/
\begin{table}[htbp]
\centering
\caption{Matched stencil population, counted by best-fit template order (which coincides with the declared scaling order for every candidate). The same candidates are evaluated under both gates; the fit label is recorded after acceptance and never used as a gate.}
\label{tab:stencil-matched}
\small
\begin{tabular}{lrrr}
\toprule
Fit order & Candidates & Stability only & Stability + Leibniz \\
\midrule
$k=0$ & 100 & 100 & 0 \\
$k=1$ & 100 & 100 & 100 \\
$k=2$ & 100 & 81 & 0 \\
Total & 300 & 281 & 100 \\
\bottomrule
\end{tabular}
\end{table}


\begin{figure}[htbp]
\centering
\includegraphics[width=\linewidth]{fig_stencil_matched.pdf}
\caption{Matched stencil experiment. (a) Normalized Leibniz defect on three grids for the three declared orders (median, with the range over 100 candidates shaded). The defect of first-order stencils falls by roughly a factor of eight per halving of $h$, while identity-like and second-order stencils keep a defect that does not decay. (b) Survivors of each order under the stability gate alone (light) and with the Leibniz gate added (solid). Data in Table~\ref{tab:stencil-matched}.}
\label{fig:stencil}
\end{figure}

The mechanism is visible in the limiting defects, independent of any fit.
For an identity-like operator the product-rule residual tends to $fg-fg-fg=-fg$, whose normalized size is about $1/3$.
For a second-derivative operator it tends to $(fg)''-f''g-fg''=2f'g'$, which on the pair $(x,x^2)$ gives exactly $1/2$ before the $10^{-12}$ denominator regularization.
For a first-derivative operator the residual tends to zero, and with $M_2=M_3=0$ it decays like $h^3$.
The Leibniz gate therefore discriminates between orders for a reason that is a theorem, not a tuning artifact.

\paragraph{Why the population must be matched.}
A tempting but unmatched comparison would contrast two different populations.
In an unconstrained random population of $\StencilBaselineTotal$ stencils (half-widths $2$ and $3$), $\StencilBaselineSurvivors$ pass the stability gate and all of them fit order~$0$, as expected, because a generic stencil acts to leading order as multiplication by $\sum_jc_j$.
In a population first projected onto the first-derivative moment conditions ($M_0=0$, $M_1=1$, and $M_2,M_3$ near zero), every one of the $\StencilLegacyTested$ candidates tested already fits order~$1$ before any gate, and the stability gate alone keeps the same $\StencilLegacySurvivors$ survivors as the stability and Leibniz gates together.
Comparing these two populations would credit the Leibniz gate with a selection that the moment normalization had already made.
The matched design holds the population fixed and varies only the gate.

\paragraph{A sampled search for false positives.}
We also searched for stencils that pass a Leibniz gate yet do not resemble any of $f$, $f'$, $f''$.
This search uses a different, pointwise defect: for $L_hf=h^{-1}\sum_jc_jf(\cdot+jh)$, the maximum over interior points and over the three pairs $(x,x^2)$, $(\sin 2\pi x,e^x)$, $(x^2,e^x)$ of
\[
  \frac{|L_h(fg)-(L_hf)g-f(L_hg)|}{|L_h(fg)|+10^{-12}},
\]
A candidate passes if the $\ell^2$ norm of its $h^{-1}$-scaled finest-grid output exceeds $10^{-6}$ for at least one test function and the three values $\Def_0,\Def_1,\Def_2$ of this defect satisfy $\Def_2<\text{threshold}$, $\Def_2<\text{ratio}\cdot\Def_1$, and $\Def_1<\text{ratio}\cdot\Def_0$; no stability gate is applied.
Among $\StencilHuntTested$ candidates ($\StencilHuntPerMode$ unconstrained and $\StencilHuntPerMode$ normalized to $M_0=0$, $M_1=1$), none passed the gate at thresholds $(0.10,0.75)$ for $(\Def_2,\text{ratio})$, so the search relaxed them to $(0.20,0.85)$.
At the relaxed thresholds $\StencilHuntPassed$ candidates pass, all from the normalized population, and none has a best template-fit error above $\StencilHuntMinF$.
This is a sampled negative result.
It is not an exhaustive classification, and a three-grid stability test does not prove bounded operator norms on a function space.

\subsection{Integration as a packaged closure of sums}
\label{subsec:integration}
The same audit applies to sums.
On the grid $x_i=ih$ in $[0,1]$, consider the cumulative left rule $I_h(f)(x_i)=h\sum_{j<i}f_j$ and the cumulative trapezoid rule $T_h(f)(x_i)=h\sum_{j<i}(f_j+f_{j+1})/2$.
They satisfy exact discrete identities:
\[
  \delta_hI_h(f)(x_i)=f_i,\qquad
  \delta_hT_h(f)(x_i)=\tfrac12(f_i+f_{i+1}),\qquad
  T_h(f)(x_i)-I_h(f)(x_i)=\tfrac h2\,(f_i-f_0),
\]
and both are exactly additive over a split of the interval.
The left rule therefore satisfies a discrete fundamental theorem with zero defect, and its measured residual (about $\IntegrationFTLeftMax$ at most), like the additivity residuals (about $\IntegrationAddMax$ at most), is floating-point rounding.
We do not fit power laws to such residuals.

The informative diagnostics, maximized over the test family $\{\sin 2\pi x,\ e^x,\ 1,\ x,\ x^2\}$ on grids with $N\in\IntegrationNList$, are the following (Table~\ref{tab:integration}, Figure~\ref{fig:integration}).
\begin{itemize}
  \item \emph{Route mismatch.} The relative mismatch $\|I_h-T_h\|/(\|T_h\|+\varepsilon_0)$ falls to $\IntegrationRMAtSmallestH$ at the finest grid, with fitted slope $\IntegrationRMExponent$ in $h$. This matches the exact identity above: for a fixed nonzero antiderivative the relative mismatch is $O(h)$, since the grid-size factors of the two norms cancel.
  \item \emph{Trapezoid fundamental theorem.} For $f\in C^1$ with $\|f'\|_\infty\le M$, the pointwise defect $|\tfrac12(f_i+f_{i+1})-f_i|$ is at most $Mh/2$. In the unweighted $\ell^2$ norm over $N=1/h$ points this is $O(h^{1/2})$, and the fitted slope is $\IntegrationFTExponent$; in the $\sqrt h$-weighted norm, which approximates the $L^2$ norm, it is $O(h)$.
  \item \emph{Truth.} Agreement between two routes does not show that either computes an integral, so we also compare each rule with the exact antiderivative. The $\sqrt h$-weighted errors decrease to $\IntegrationLeftTruthFinest$ (left) and $\IntegrationTrapTruthFinest$ (trapezoid) at $N=1024$, consistent with the classical $O(h)$ and $O(h^2)$ error orders for $C^1$ and $C^2$ integrands.
\end{itemize}

% Auto-generated by scripts/export_results_tex.py from notes/math_review/
\begin{table}[htbp]
\centering
\caption{Integration diagnostics, each maximized over the test family. RM is the relative left/trapezoid route mismatch; $\mathrm{FT}_T$ is the trapezoid fundamental-theorem defect in the raw and $\sqrt h$-weighted $\ell^2$ norms; the last two columns are $\sqrt h$-weighted errors against the exact antiderivative.}
\label{tab:integration}
\small
\begin{tabular}{rrrrrr}
\toprule
$N$ & RM & $\mathrm{FT}_T$ raw & $\mathrm{FT}_T$ weighted & Left error & Trapezoid error \\
\midrule
128 & $0.0142$ & $0.196$ & $0.0174$ & $3.42\times 10^{-3}$ & $3.91\times 10^{-5}$ \\
256 & $7.09\times 10^{-3}$ & $0.139$ & $8.68\times 10^{-3}$ & $1.71\times 10^{-3}$ & $9.79\times 10^{-6}$ \\
512 & $3.54\times 10^{-3}$ & $0.0982$ & $4.34\times 10^{-3}$ & $8.52\times 10^{-4}$ & $2.45\times 10^{-6}$ \\
1024 & $1.77\times 10^{-3}$ & $0.0694$ & $2.17\times 10^{-3}$ & $4.25\times 10^{-4}$ & $6.12\times 10^{-7}$ \\
\bottomrule
\end{tabular}
\end{table}


\begin{figure}[htbp]
\centering
\includegraphics[width=0.92\linewidth]{fig_integration.pdf}
\caption{Integration diagnostics under refinement (log--log). Annotated slopes are the theoretical orders. The left-rule fundamental-theorem residual (grey) is rounding error around an exact identity and is shown only for contrast. Data in Table~\ref{tab:integration}.}
\label{fig:integration}
\end{figure}
